\documentclass[a4paper, 11pt]{article}

% -- Coments
\usepackage{verbatim}

% ----- Fonts -----
% -- Fuente --
% \usepackage{fontspec}
% \setmonofont{JetBrainsMono Nerd Font}

\usepackage{babel}
\babelprovide[import, main]{spanish}

% -- Color --
\usepackage{xcolor}
%\definecolor{azul}{RGB}{00,33,99}
\definecolor{azul}{RGB}{35,72,180}

% -- Page Margin --
\usepackage[margin=1in]{geometry}

% -- Espaciados --
\newcommand{\Pspace}{0.5cm}
\newcommand{\Aspace}{0.2cm}

% -- Columnas --
\usepackage{multicol}

% -- Imagenes --
\usepackage{graphicx}
\usepackage{float}

% -- Matemáticas --
\usepackage{amsmath, amssymb}

% -- Gráficas --
\usepackage{pgfplots}
\pgfplotsset{compat=1.18}


\title
{
    Contexto Regional 2026-2 \\
    Algarada femenina en la posrevolución sonorense - Zúñiga Elizalde, Mercedes \\

    Profesor José Moreno Vega
}

\begin{document}

\begin{otherlanguage}{spanish}
\maketitle
\end{otherlanguage}

\begin{center}
    \begin{tabular}{r|l}
        \textbf{Expediente} & \textbf{Nombre} \\ \hline
        219208106 & Bórquez Guerrero Angel Fernando \\
    \end{tabular}
\end{center}

{\Large 1. ¿Cuál es el argumento central de la autora?} \\
{\color{azul}
    La autora dice que en la posrevolución sonorense, sobre todo en los años treinta, cuarenta y cincuenta, las mujeres tuvieron una participación política activa aunque todavía no tenían derechos ciudadanos, y que lo hicieron desde el magisterio, los sindicatos y la prensa. Los gobiernos impulsaron su organización, pero siempre dentro de los cauces de sus propios proyectos políticos. Para mostrarlo sigue a cinco mujeres (Enriqueta de Parodi, Catalina Acosta, Dolores Enedina Cuiltre, María Jesús Guirado y Francisca Córdova), que construyeron un liderazgo "a contracorriente" haciendo alianzas con políticos y gobernantes, hasta llegar a ser diputadas cuando se aprobó el voto en 1953.
}
\vspace{1cm}

{\Large 2. ¿Qué evidencias utiliza para respaldarlo?} \par
{\color{azul}
Aunque dice que hay poca documentación, usa sobre todo documentos del Archivo General del Estado de Sonora, como los fondos de Oficialía Mayor, Gobernación, Cultos y Templos y Educación, donde hay cartas de los sindicatos a los gobernadores, informes de gobierno y oficios. También usa la prensa de la época, como \textit{El Pueblo}, \textit{El Tiempo} y \textit{Ahora}, y los artículos que escribieron Parodi y Acosta, donde se ve cómo pensaban sobre el lugar de la mujer. Presenta casos concretos, como el de Guirado, secretaria general de la CTS y primera mujer diputada local en 1955, y el de Cuiltre, que fue despedida por apoyar una huelga. Por último, se apoya en estudios de otros autores como Cejudo y Vaughan.
}
\vspace{1cm}

{\Large 3. ¿Cómo se relaciona esta lectura con la lectura anterior?} \par
{\color{azul}
La lectura 8 hablaba de cómo Obregón, Calles y su red familiar controlaban Sonora, y aquí aparecen esos mismos gobiernos, como el de Rodolfo Elías Calles con su campaña desfanatizadora, vistos desde las maestras y sindicalistas que tuvieron que moverse dentro de ese poder. Además, regresa el tema de la lectura 7: las mujeres que antes quedaban en el silencio ahora aparecen como protagonistas.
}

\end{document}
